% Part: model-theory
% Chapter: basics
% Section: isomorphism

\documentclass[../../../include/open-logic-section]{subfiles}

\begin{document}

\olfileid{mod}{bas}{iso}
\olsection{ساختارهای همریخت}

!!{structure}s مرتبهٔ اول از دو جهت می‌توانند همانند باشند. یک جهتِ
همانندی آن است که !!{sentence}s یکسانی را صادق سازند. چنین
!!{structure}sی را \emph{هم‌ارز ابتدایی} می‌نامیم. اما ساختارها
می‌توانند بسیار متفاوت باشند و بااین‌حال همان !!{sentence}s را صادق
سازند---برای نمونه، یکی ممکن است !!{enumerable} باشد و دیگری نباشد.
علت این است که ویژگی‌های فراوانی از !!a{structure} را نمی‌توان در
زبان‌های مرتبهٔ اول بیان کرد، خواه به این سبب که زبان به‌اندازهٔ کافی
غنی نیست و خواه به سبب محدودیت‌های بنیادی منطق مرتبهٔ اول، مانند قضیهٔ
لوونهایم--اسکولم. پس جهتِ دیگر و سخت‌گیرانه‌تری که در آن
!!{structure}s می‌توانند همانند باشند این است که از بنیاد یکسان باشند؛
یعنی تنها در اشیای تشکیل‌دهنده‌شان تفاوت کنند، نه در ویژگی‌های
ساختاری‌شان. راهی برای دقیق‌کردن این معنا مفهوم \emph{همریختی} است.

\begin{defn}
\ollabel{defn:elem-equiv}
دو !!{structure}s~$\Struct{M}$ و~$\Struct M'$ را برای
!!{language} یکسان~$\Lang{L}$ در نظر بگیرید. می‌گوییم $\Struct{M}$
\emph{با}~$\Struct M'$ \emph{هم‌ارز ابتدایی است}، و می‌نویسیم
$\Struct{M} \equiv \Struct M'$، هرگاه و تنها هرگاه برای هر
!!{sentence}~$!A$ از~$\Lang{L}$، $\Sat{M}{!A}$ هرگاه و تنها هرگاه
$\Sat{M'}{!A}$.
\end{defn}

\begin{defn}
\ollabel{defn:isomorphism} دو !!{structure}s~$\Struct{M}$ و
$\Struct M'$ را برای !!{language} یکسان~$\Lang L$ در نظر بگیرید.
می‌گوییم $\Struct{M}$ \emph{با}~$\Struct M'$ \emph{همریخت است}، و
می‌نویسیم $\Struct{M} \simeq \Struct M'$، هرگاه و تنها هرگاه تابعی
$h \colon \Domain{M} \to \Domain{M'}$ با شرایط زیر وجود داشته باشد:
\begin{enumerate}
\item $h$ !!{injective} است: اگر $h(x) = h(y)$، آنگاه $x = y$؛
\item $h$ !!{surjective} است: برای هر $y \in \Domain{M'}$،
  $x \in \Domain{M}$ای وجود دارد که $h(x) = y$؛
\item \ollabel{defn:iso-const}برای هر !!{constant}~$c$:
  $h(\Assign{c}{M}) = \Assign{c}{M'}$؛
\item \ollabel{defn:iso-pred}برای هر !!{predicate}~$n$موضعی $P$:
  \[
  \tuple{a_1, \dots, a_n}\in \Assign{P}{M} \quad\text{هرگاه و تنها هرگاه}\quad
  \tuple{h(a_1), \dots, h(a_n)} \in \Assign{P}{M'};
  \]
\item \ollabel{defn:iso-func}برای هر !!{function}~$n$موضعی $f$:
  \[
  h(\Assign{f}{M}(a_1, \dots, a_n)) =
  \Assign{f}{M'}(h(a_1), \dots, h(a_n)).
  \]
\end{enumerate}
\end{defn}

\begin{thm}
\ollabel{thm:isom}
اگر $\Struct{M} \iso \Struct M'$، آنگاه $\Struct{M} \elemequiv
\Struct{M'}$.
\end{thm}

\begin{proof}
$h$ را همریختی‌ای از~$\Struct{M}$ به روی~$\Struct M'$ بگیرید. برای هر
گمارش~$s$، $h \circ s$ ترکیب $h$ و~$s$ است؛ یعنی گمارشی در
$\Struct{M'}$ که $(h \circ s)(x) = h(s(x))$. با استقرا بر $t$ و~$!A$
می‌توان ادعاهای قوی‌تر زیر را اثبات کرد:
\begin{enumerate}
  \item[a.] $h(\Value{t}{M}[s]) = \Value{t}{M'}[h\circ s]$.
  \item[b.] $\Sat{M}{!A}[s]$ هرگاه و تنها هرگاه
  $\Sat{M'}{!A}[h \circ s]$.
\end{enumerate}
ادعای نخست با استقرا بر پیچیدگی~$t$ اثبات می‌شود.
\begin{enumerate}
\item اگر $t \ident c$، آنگاه $\Value{c}{M}[s] = \Assign{c}{M}$ و
  $\Value{c}{M'}[h \circ s] = \Assign{c}{M'}$. پس
  $h(\Value{t}{M}[s]) = h(\Assign{c}{M}) = \Assign{c}{M'}$ (بنا بر
  \olref{defn:iso-const} از \olref{defn:isomorphism}) $=
  \Value{t}{M'}[h \circ s]$.
\item اگر $t \ident x$، آنگاه $\Value{x}{M}[s] = s(x)$ و
  $\Value{x}{M'}[h \circ s] = h(s(x))$. پس
  $h(\Value{x}{M}[s]) = h(s(x)) = \Value{x}{M'}[h \circ s]$.
\item اگر $t \ident f(t_1, \dots, t_n)$، آنگاه
  \begin{align*}
    \Value{t}{M}[s] & = \Assign{f}{M}(\Value{t_1}{M}[s], \dots, \Value{t_n}{M}[s]) \quad\text{و}\\
  \Value{t}{M'}[h \circ s] & = \Assign{f}{M'}(\Value{t_1}{M'}[h \circ
    s], \dots, \Value{t_n}{M'}[h \circ s]).
  \end{align*}
  فرض استقرا این است که برای هر $i$،
  $h(\Value{t_i}{M}[s]) = \Value{t_i}{M'}[h\circ s]$. بنابراین،
  \begin{align}
    h(\Value{t}{M}[s])
    & = h(\Assign{f}{M}(\Value{t_1}{M}[s], \dots, \Value{t_n}{M}[s])) \notag\\
    & = \Assign{f}{M'}(h(\Value{t_1}{M}[s]), \dots,
    h(\Value{t_n}{M}[s])) \ollabel{iso-1}\\
    & = \Assign{f}{M'}(\Value{t_1}{M'}[h \circ s], \dots,
    \Value{t_n}{M'}[h \circ s]) \ollabel{iso-2}\\
    & = \Value{t}{M'}[h\circ s] \notag
  \end{align}
  در اینجا، \olref{iso-1} از \olref{defn:iso-func} در
  \olref{defn:isomorphism} و \olref{iso-2} از فرض استقرا نتیجه می‌شود.
\end{enumerate}
بخش (b) به‌عنوان تمرین واگذاشته می‌شود.

اگر $!A$ یک جمله باشد، گمارش‌های $s$ و~$h \circ s$ بی‌اثرند و
$\Sat{M}{!A}$ هرگاه و تنها هرگاه $\Sat{M'}{!A}$.
\end{proof}

\begin{prob}
برهان بخش (b) از \olref[mod][bas][iso]{thm:isom} را با جزئیات کامل
کنید. حتماً مشخص کنید هر یک از پنج خاصیتی که همریختی‌ها را در
\olref[mod][bas][iso]{defn:isomorphism} مشخصه‌سازی می‌کنند در کجا به
کار می‌رود.
\end{prob}

\begin{defn}
یک \emph{خودریختی} ساختار~$\Struct{M}$ همریختی‌ای از~$\Struct{M}$ به
روی خودش است.
\end{defn}

\begin{prob}
نشان دهید برای هر ساختار~$\Struct{M}$، اگر $X$ زیرمجموعه‌ای تعریف‌پذیر
از~$\Struct{M}$ و $h$ خودریختی‌ای از~$\Struct{M}$ باشد، آنگاه
$X = \Setabs{h(x)}{x \in X}$ (یعنی $X$ تحت~$h$ ثابت می‌ماند).
\end{prob}


\end{document}
